\begin{enumerate}
	\item 
	\begin{enumerate}
		\item On calcule les coordonnées de $\vect{\text{AB}}$ et de $\vect{\text{AC}}$.
		
		$\vect{\text{AB}}\left(
		\begin{array}{c}
		-5\\
		2\\
		-5
		\end{array}
		\right)$
et
		$\vect{\text{AC}}\left(
		\begin{array}{c}
		-4\\
		5\\
		2
		\end{array}
		\right)$
		
		On constate que $\dfrac{-5}{2}\neq \dfrac{-4}{5}$. Comme les coordonnées des vecteurs ne sont pas proportionnelles, 	$\vect{\text{AB}}$ et $\vect{\text{AC}}$ ne sont pas colinéaires et A, B et C ne sont pas alignés.
		\item
D appartient au plan (ABC) si et seulement si les coordonnées de D vérifient l'équation cartésienne du plan (ABC).

$ 29\times (-3) + 30 \times (-4) - 17 \times 6 = -309 \neq 35$. Les coordonnées de D ne vérifient pas  l'équation cartésienne du plan (ABC), donc les points A, B, C, D ne sont pas coplanaires.
	\end{enumerate}
	\item 
	\begin{enumerate}
		\item Les coordonnées du milieu de [AB] sont :
$\dfrac{4+(-1)}{2}~;~\dfrac{-1+1}{2}~;~\dfrac{3+(-2)}{2}$ soit $\dfrac32~;~0~;~\dfrac12$.
	
		\item Un vecteur normal au plan $P_1$ médiateur du segment [AB] est le vecteur $\vect{\text{BA}}\left(
		\begin{array}{c}
		5\\
		-2\\
		5
		\end{array}
		\right)$ donc une équation de $P_1$ est $5x - 2y + 5z = k$ avec $k$ tel que les coordonnées du milieu de [AB] vérifient l'équation c'est-à dire si 

$5\times\dfrac{3}{2} - 2\times 0 + 5\times\dfrac{1}{2} =  k =10$

Ainsi $5x - 2y + 5z = 10$ est une équation cartésienne de $P_1$.
	\end{enumerate}
	\item
	\begin{enumerate}
	\item \ {}

	$\begin{array}{ccl}\text{MC}^2 &=& (x - 0)^2 + (y - 4)^2 + (z - 5)^2\\ &=& x^2 + y^2 + z^2 -8y -10z + 16 + 25 \\ &=& x^2 + y^2 + z^2 -8y -10z + 41\end{array}$

	$\begin{array}{ccl}\text{MD}^2 &=&  (x - (-3))^2 + (y - (-4))^2 + (z - 6)^2\\ &=&  x^2 + y^2 + z^2 + 6x + 8y -12z + 9 + 16 + 36 \\&=&  x^2 + y^2 + z^2 + 6x + 8y - 12z + 61\end{array}$
	
	M appartient à $P_2$ est équivalent à $\text{MC}^2 = \text{MD}^2$.
	
	$
	\begin{array}{ccl}
	\text{MC}^2 = \text{MD}^2 & \iff & \text{MC}^2 - \text{MD}^2 = 0\\
	& \iff & x^2 + y^2 + z^2 -8y -10z + 41 \\&&- x^2 - y^2 - z^2 - 6x - 8y + 12z - 61 = 0\\
	& \iff & - 6x - 16y + 2z - 20 = 0\\
	& \iff & - 6x - 16y + 2z  = 20\\
	& \iff & - 3x - 8y + z  = 10\\
	\end{array}
	$
	
	On obtient l'équation cartésienne de $P_2$ : $- 3x - 8y + z  = 10$.
	\item Un vecteur normal à $P_1$ est le vecteur de coordonnées $\left(
	\begin{array}{c}
	5\\
	-2\\
	5
	\end{array}
	\right)$.
	
	Un vecteur normal à $P_2$ est le vecteur de coordonnées $\left(
	\begin{array}{c}
	-3\\
	-2\\
	5
	\end{array}
	\right)$
	
	On remarque que $\dfrac{5}{-2} \neq \dfrac{-3}{-2}$
	
	Ces deux vecteurs ne sont pas colinéaires car leur coordonnées ne sont pas proportionnelles.
		
	Comme les vecteurs normaux à $P_1$ et  $P_2$ ne sont pas colinéaires, alors les plans $P_1$ et  $P_2$ sont sécants.
	\end{enumerate}
	\item $\forall t \in \R\quad 5\times(-2-\np{1.9}t) - 2\times t + 5 \times (4+\np{2.3}t) = -10 -\np{9.5}t- 2t + 20 +  \np{11.5}t = 10 $
	
	Tout point de $\Delta$ appartient à $P_1$.
	
	$\forall t \in \R\quad -3\times(-2-\np{1.9}t) - 8\times t + (4+\np{2.3}t) = 6 +\np{5.7}t- 8t + 4 +  \np{2.3}t = 10 $
	
	Tout point de $\Delta$ appartient à $P_2$.
	
	$\Delta$ est donc la droite d'intersection de   $P_1$ et $P_2$.
	\item D'après la représentation paramétrique de $\Delta$, un vecteur directeur de $\Delta$ est le vecteur $\vect{\delta}$ de coordonnées $\left(
		\begin{array}{c}
		19\\
		-10\\
		23
		\end{array}
		\right)$.
		
		Le plan  $P_3$ admet un vecteur normal $\vect{n_3}$ de coordonnées $\left(
		\begin{array}{c}
		8\\
		-10\\
		-4
		\end{array}
		\right)$.
		
La droite $\Delta$ et le plan $P_3$ ne sont pas sécants si et seulement si un vecteur directeur de $\Delta$ est orthogonal à un vecteur normal au plan $P_3$.
		
On calcule le produit scalaire des deux vecteurs précités :
\renewcommand\arraystretch{1}

$\vect{\delta} \cdot \vect{n_3} = 19 \times 8 + (-10) \times (-10) + 23 \times (-4) = 160 \neq 0$

		
Le produit scalaire n'est pas nul, les deux vecteurs ne sont pas orthogonaux, $\Delta$ et $P_3$ sont donc sécants.
	\item On considère le point H équidistant des points A, B, C, D. Comme AH~=~BH, H appartient à $P_1$. Comme DH~=~CH, H appartient à $P_2$. Ainsi H appartient à l'intersection des deux plans, la droite $\Delta$. Comme AH~=~CH, H appartient au plan $P_3$. Ainsi H est l'intersection de  $\Delta$ et de  $P_3$.
	
	\emph{Le point équidistant des sommets d'un tétraèdre est l'intersection des plans médiateurs de ses arêtes.}
	
On peut déterminer les coordonnées de H.
	
Il suffit de résoudre l'équation :
	
$8(-2-\np{1.9}t) -10t -4(4+\np{2.3}t) = -15$


\end{enumerate}
